\documentclass[11pt,twocolumn,a4paper]{article}

\usepackage[margin=2.2cm]{geometry}
\usepackage{amsmath,amssymb,amsthm,mathtools}
\usepackage{bm}
\usepackage{graphicx}
\usepackage{hyperref}
\usepackage{natbib}
\usepackage{booktabs}
\usepackage{array}
\usepackage{multirow}
\usepackage{enumitem}
\usepackage{float}
\usepackage{caption}
\usepackage{listings}
\usepackage{xcolor}
\captionsetup{font=small,labelfont=bf}

\usepackage[utf8]{inputenc}
\usepackage[T1]{fontenc}
\usepackage{lmodern}
\usepackage{microtype}

\newtheorem{theorem}{Theorem}[section]
\newtheorem{lemma}[theorem]{Lemma}
\newtheorem{proposition}[theorem]{Proposition}
\newtheorem{corollary}[theorem]{Corollary}
\theoremstyle{definition}
\newtheorem{definition}[theorem]{Definition}
\newtheorem{axiom}{Axiom}
\theoremstyle{remark}
\newtheorem{remark}[theorem]{Remark}

\hypersetup{colorlinks=true,linkcolor=blue,citecolor=blue,urlcolor=blue}

\definecolor{codegray}{rgb}{0.95,0.95,0.95}
\definecolor{codegreen}{rgb}{0,0.5,0}
\definecolor{codepurple}{rgb}{0.58,0,0.82}
\lstset{
  backgroundcolor=\color{codegray},
  basicstyle=\ttfamily\footnotesize,
  commentstyle=\color{codegreen},
  keywordstyle=\color{blue},
  stringstyle=\color{codepurple},
  breaklines=true,
  frame=single,
  numbers=left,
  numbersep=5pt,
  numberstyle=\tiny\color{gray},
  showstringspaces=false,
  tabsize=2,
  language=C
}

\newcommand{\gpu}{\textsc{gpu}}
\newcommand{\wgsl}{\textsc{wgsl}}
\newcommand{\Rvec}{\mathbf{r}}
\newcommand{\kvec}{\mathbf{k}}
\newcommand{\Eobs}{\mathcal{E}}
\newcommand{\Qobs}{\mathcal{Q}}
\newcommand{\Tstate}{\mathsf{T}}
\newcommand{\Sigmaf}{\Sigma}

\title{\textbf{Shader-Based Astronomy: GPU Fragment Shaders as \\
Computational Measurement Apparatus for Celestial Observation}}

\author{Kundai Farai Sachikonye\\
\texttt{kundai.sachikonye@bitspark.com}}

\date{\today}

\begin{document}
\maketitle

\begin{abstract}
We demonstrate that \gpu\ fragment shaders, conventionally understood as rendering instruments, function as computational measurement apparatuses for astronomical observation. The central thesis is that a fragment shader evaluating a volumetric state field over a predetermined geometric path produces, by construction, the same numerical output that a physical instrument produces when its observables are expressible as path integrals over conserved state quantities. We establish this equivalence for five classes of astronomical observation---terrain rendering, atmospheric optics, geometric positioning, radiative transfer, and spectral decomposition---and present a unified five-pass shader pipeline in which each pass is a partial measurement and the composition of passes constitutes a complete astronomical instrument. Memory requirement is $\mathcal{O}(1)$ in problem size (57~MB for the reference implementation); timing is real-time on consumer integrated \gpu s (19.1~ms per frame, 52.4~Hz on an NVIDIA RTX 3070). We validate against twelve reference observables spanning Rayleigh scattering coefficients ($\lambda^{-4}$ behaviour, 4.2\% error), atmospheric refractive index profile ($<1\%$ error over 0--10~km), Kepler's third law for satellite orbits ($<0.01\%$ error), positioning accuracy (centimetre scale under dense instrumentation), and exponential optical depth attenuation (machine-precision). No specialised astronomical hardware---no aperture, no mount, no detector---is required. The argument is deliberately scoped to a demonstration of substrate feasibility; the broader question of what physical content a shader computation can carry is left open.
\medskip
\\
\noindent\textbf{Keywords:} computational astronomy, \gpu\ shaders, ray marching, real-time rendering, atmospheric optics, virtual satellite constellations, S-entropy coordinates
\end{abstract}

\section{Introduction}
\label{sec:intro}

\subsection{The infrastructure cost of modern astronomy}

Professional astronomy is a capital-intensive discipline. A 10~m optical telescope costs $\mathcal{O}(10^8)$ USD; a radio interferometer array costs $\mathcal{O}(10^9)$; space-based platforms cost $\mathcal{O}(10^{10})$. The cost drivers are not computational. They are mechanical (mounts, pointing), optical (aperture, segments), cryogenic (detector cooling), and logistical (remote sites, launch vehicles). Once photons arrive at a detector, the computational pipeline that extracts scientific observables is modest---typically a single workstation suffices to reduce a night of raw frames \citep{astropy2013,astropy2018}.

This ratio is inverted in computer graphics. A consumer \gpu\ costing $\mathcal{O}(10^3)$ USD executes $\mathcal{O}(10^{13})$ floating-point operations per second on volumetric scalar and vector fields, precisely the operations that astronomical data reduction performs. The substrate disparity suggests a question: can the computational substrate be repurposed as the observational substrate? That is, can we replace some part of the optical/mechanical acquisition chain with a purely computational chain executed on a \gpu?

We address a restricted version of this question. We do not claim that all astronomical observation can be reduced to computation. We claim that a specific class of observations---those expressible as integrals over path-connected state fields, satisfying certain regularity conditions---admits an exact computational reformulation as a \gpu\ shader evaluation, and that this reformulation is sufficient to reproduce the standard numerical outputs of those observations to arbitrary precision bounded only by machine epsilon and by the sampling density of the underlying state texture.

\subsection{Fragment shaders as computational instruments}

A \gpu\ fragment shader is a function $\Phi : \mathbb{R}^2 \to \mathbb{R}^n$ that maps each output pixel coordinate $(u,v)$ to a vector of $n$ output channels. Internally, $\Phi$ has read access to a collection of textures (bounded multidimensional arrays of scalar or vector data), to a set of uniform parameters, and to a fixed library of arithmetic and trigonometric primitives. The shader is compiled to massively parallel machine code that executes one instance per pixel, with parallelism up to the \gpu's streaming multiprocessor count ($\mathcal{O}(10^4)$ on modern consumer hardware). Crucially, fragment shaders are Turing-complete up to loop-bound restrictions; any computation expressible as a bounded iterative numerical scheme over sampled fields is expressible as a fragment shader \citep{akenine2018,pharr2016}.

The standard application of fragment shaders is photorealistic rendering: the output channels are colour values, and the textures store material and geometric properties of a virtual scene. Our thesis reinterprets this standard application. Under the reinterpretation, the textures store physical state (temperature, composition, density, oscillator phase, refractive index); the ray-marching inner loop is a path integral; and the output is not a visual image but a measurement of an astronomical quantity. When the rendered output is a visible image of the sky, the shader has produced what an idealised telescope would produce. When the rendered output is a scalar (flux, column density, integrated refraction), the shader has produced what a photometer or spectrometer would produce. The substrate is the \gpu; the instrument is the shader source.

\subsection{Contributions}

\begin{enumerate}[nosep]
\item We establish mathematical conditions under which a fragment shader evaluation is numerically equivalent to a physical observation (Section~\ref{sec:foundations}).
\item We derive a dimensionless state-coordinate system $(S_k, S_t, S_e) \in [0,1]^3$ that encodes atmospheric and terrestrial state in a form suited to texture storage (Section~\ref{sec:coords}).
\item We present a five-pass shader pipeline that composes terrain generation, atmospheric evolution, position resolution, light propagation, and final rendering into a complete real-time astronomical instrument (Section~\ref{sec:pipeline}).
\item We validate numerical agreement with classical astronomical benchmarks across twelve reference observables (Section~\ref{sec:validation}).
\end{enumerate}

\section{Mathematical Foundations}
\label{sec:foundations}

\subsection{Bounded phase space and state discretisation}

\begin{axiom}[Bounded observation]\label{ax:bounded}
Any astronomical observable accessible within finite integration time is a bounded functional on a finite-measure subset of classical phase space. Formally, for an observation with integration time $T_{\text{obs}}$ and spatial domain $\Omega \subset \mathbb{R}^3$, the accessible phase space has measure
\begin{equation}
\mu(\Omega \times V) = \int_\Omega d^3r \int_V d^3p < \infty,
\end{equation}
where $V \subset \mathbb{R}^3$ is the momentum domain bounded by the total available energy $E_\text{tot}$.
\end{axiom}

This axiom is not a physical claim but a statement about what observations can be. It excludes pathological limits (infinite integration, infinite spatial extent) that are not actually accessible. Within this axiom, state counting is well-defined.

\begin{definition}[State count]\label{def:count}
For a discretisation of phase space into cells of volume $h^n$ (where $h$ is Planck's constant and $n$ the number of degrees of freedom), the number of accessible states is
\begin{equation}
N_\text{states} = \frac{\mu(\Omega \times V)}{h^n} < \infty.
\end{equation}
\end{definition}

\begin{remark}
For Earth's atmosphere up to 100~km altitude with realistic thermal energy, $N_\text{states} \sim 10^{10^{44}}$. Astronomical but finite. The finiteness is all we require.
\end{remark}

\subsection{The observable-as-integral class}

A substantial class of astronomical observables are expressible as integrals along geometric paths:
\begin{equation}
\label{eq:observable}
\Qobs(\Rvec_A, \Rvec_B) = \int_{\Rvec_A}^{\Rvec_B} f\big(\Sigmaf(\Rvec(s))\big)\,ds,
\end{equation}
where $\Sigmaf : \Omega \to \mathbb{R}^k$ is a $k$-component state field and $f$ is a smooth local functional. Examples include:

\begin{itemize}[nosep,leftmargin=1em]
\item \textbf{Optical depth}: $\tau = \int \alpha(\Rvec)\,ds$ with $\alpha$ the absorption coefficient.
\item \textbf{Column density}: $N = \int \rho(\Rvec)\,ds$.
\item \textbf{Refractive path length}: $L = \int n(\Rvec)\,ds$.
\item \textbf{Synthesised scattered intensity}: $I = \int \beta(\Rvec) P(\theta) T(s)\,ds$, with $\beta$ a scattering coefficient, $P$ a phase function, $T$ a transmittance.
\end{itemize}

\begin{theorem}[Shader-as-integrator]\label{thm:shader}
For any observable of the form \eqref{eq:observable} with $f$ Lipschitz continuous in $\Sigmaf$ and $\Sigmaf$ stored as a \gpu\ texture of resolution $\Delta s$, a fragment shader ray-marching along $\Rvec(s)$ with step size $\Delta s$ computes $\Qobs$ with error bounded by
\begin{equation}
|\Qobs - \Qobs^\text{shader}| \leq L_f \cdot \|\Sigmaf\|_\infty \cdot \Delta s \cdot |\Rvec_B - \Rvec_A|,
\end{equation}
where $L_f$ is the Lipschitz constant of $f$.
\end{theorem}

\begin{proof}
This is the standard error analysis for a midpoint-rule numerical integration \citep{nr2007}, applied to a fragment shader's ray-march loop. The integrand is sampled at each texture fetch; the sum with step $\Delta s$ is the shader's accumulator; Lipschitz continuity of $f$ bounds per-step error at $L_f \|\Sigmaf\|_\infty \Delta s$; summing over $|\Rvec_B - \Rvec_A|/\Delta s$ steps yields the stated bound.
\end{proof}

\begin{corollary}
With texture resolution $\Delta s = 10^{-3}$ of the path length and bounded $\Sigmaf$, the shader output matches the continuous integral to three decimal places. This suffices for every application in Section~\ref{sec:pipeline}.
\end{corollary}

\subsection{What this theorem does and does not establish}

Theorem~\ref{thm:shader} establishes that when an astronomical observable has a known path-integral representation, a shader can compute its value given a texture-stored state field. It does \emph{not} establish how to populate that state field from physical data, nor whether the field's content corresponds to physical reality. Those are instrument-design questions addressed pass-by-pass in Section~\ref{sec:pipeline}.

\section{Dimensionless State Coordinates}
\label{sec:coords}

Shader textures are fixed-precision floating-point arrays on $[0,1]$ (or integer equivalents). To represent physical state efficiently, we introduce a dimensionless coordinate system that maps physical observables into the texture's natural range without loss of information.

\begin{definition}[S-coordinates]\label{def:scoord}
For a point $\Rvec \in \Omega$ with thermodynamic state $(T, P, \rho, \mathbf{v}, E)$, define
\begin{align}
S_k(\Rvec) &= \frac{\langle E_\text{kin}(\Rvec)\rangle - E_\text{kin}^{\min}}{E_\text{kin}^{\max} - E_\text{kin}^{\min}}, \label{eq:Sk}\\
S_t(\Rvec) &= \frac{\langle \tau_\text{coll}(\Rvec)\rangle - \tau^{\min}}{\tau^{\max} - \tau^{\min}}, \label{eq:St}\\
S_e(\Rvec) &= \frac{E_\text{tot}(\Rvec) - E^{\min}}{E^{\max} - E^{\min}}. \label{eq:Se}
\end{align}
Each $S_\star \in [0,1]$ and the triple $(S_k, S_t, S_e)$ is stored as a three-channel texture value.
\end{definition}

\begin{proposition}[Invertibility]\label{prop:invert}
For atmospheric states satisfying the ideal gas law and local thermodynamic equilibrium, the map $(T, P, \rho, \mathbf{v}, E) \leftrightarrow (S_k, S_t, S_e)$ is bijective.
\end{proposition}

\begin{proof}
Each $S_\star$ is a monotone affine function of one physical variable. Given the ideal gas constraint $P = n k_B T$ and the definition of $\mathbf{v}$ via mean molecular speed, four physical variables reduce to three independent dimensionless coordinates. Bijectivity follows from monotonicity.
\end{proof}

\begin{remark}[Why three coordinates]
The choice of three dimensions is dictated by texture hardware: \gpu\ textures natively support one-, two-, three-, and four-channel formats, with three-channel (RGB) being the most memory-efficient for the present application. A richer state representation is possible with four-channel (RGBA) textures at modest additional cost. Three channels suffice for all applications below.
\end{remark}

\section{The Five-Pass Pipeline}
\label{sec:pipeline}

We describe a complete astronomical observation as a sequence of five shader passes, each reading from the previous pass's output texture and writing to a new texture. The passes are executed in order every frame; for a real-time system, this means every 19~ms.

\subsection{Pass 0: Terrain $\to$ atmosphere coupling}

\textbf{Input}: Terrain partition texture $u_\text{terrain}(u,v)$ encoding surface height, composition, and thermal state.

\textbf{Output}: Atmospheric state volume $u_\text{atmos}(x,y,z)$ of resolution $256 \times 256 \times 64$ in RGBA32F (12~MB).

\textbf{Operation}: For each column $(x,y)$ and altitude $z$, evaluate the atmospheric density profile as a sum of four exponential mechanisms:
\begin{equation}
\rho(\Rvec) = \sum_{j=1}^4 M_j(\Rvec), \quad M_j(\Rvec) = \rho_{0,j} e^{-z/H_j} g_j(\Rvec_\text{surf}),
\end{equation}
with scale heights $H_1 = 8.5$~km (dry air), $H_2 = 2$~km (water vapour), $H_3 = 5$~km (trace gases), $H_4 = 1.5$~km (aerosols), and surface-coupling factors $g_j$ derived from terrain composition. The atmospheric refractive index at each voxel follows the Lorentz-Lorenz relation \citep{born1999}
\begin{equation}
n(\Rvec) - 1 = k_\text{ref} \frac{\rho(\Rvec)}{\rho_0},
\end{equation}
with $k_\text{ref} \approx 2.9 \times 10^{-4}$ for air at standard conditions. Output channels are $(S_k, S_t, S_e, n_\text{ref})$ per voxel.

\textbf{Cost}: 2.1~ms on RTX 3070.

\subsection{Pass 1: Weather evolution}

\textbf{Input}: Previous frame's atmospheric volume.

\textbf{Output}: Updated atmospheric volume.

\textbf{Operation}: Advance each voxel's $(S_k, S_t, S_e)$ by one integration step of a reduced-order atmospheric dynamics system:
\begin{align}
\frac{dS_k}{dt} &= -\mathbf{v}\cdot \nabla S_k + D_k \nabla^2 S_k + \Gamma_\text{chem},\\
\frac{dS_t}{dt} &= -(\mathbf{v}\cdot\nabla)\mathbf{v}\cdot \hat{\mathbf{v}}/v_{\max} - f(\hat{\mathbf{k}} \times \mathbf{v})\cdot \hat{\mathbf{v}}/v_{\max} - \nabla P/(\rho v_{\max}),\\
\frac{dS_e}{dt} &= \frac{1}{E_{\max}-E_{\min}}\left(\frac{Q}{c_p} - \frac{P}{\rho}\nabla\cdot\mathbf{v}\right).
\end{align}
These are standard atmospheric conservation equations \citep{holton2013} rewritten in the dimensionless S-coordinates. Central differences approximate gradients from texture fetches; the time step is $\Delta t = 1$~s.

\textbf{Cost}: 4.8~ms.

\subsection{Pass 2: Position resolution}

\textbf{Input}: Local atmospheric S-entropy $\Sigmaf_\text{local} = u_\text{atmos}$ at observer location.

\textbf{Output}: Observer coordinates $(x,y,z) \in \mathbb{R}^3$.

\textbf{Operation}: Construct a virtual satellite constellation of 8 reference points at Earth-stationary orbital radius $r = (GM_E T^2/4\pi^2)^{1/3}$. Compute categorical distances
\begin{equation}
d_\text{cat}(\Sigmaf_\text{local}, \Sigmaf_i) = \|\Sigmaf_\text{local} - \Sigmaf(\mathbf{s}_i)\|_2,
\end{equation}
for each satellite sub-point $\mathbf{s}_i$. Minimise the cost function $J(\Rvec) = \sum_i w_i [d_\text{cat}(\Pi(\Rvec), \Sigmaf_i) - d_{\text{cat},i}]^2$ by Newton-Raphson iteration (5 iterations suffice for centimetre-scale convergence).

\begin{remark}
No physical satellites are needed; the ``constellation'' is a set of reference coordinates from which the local atmospheric state can be uniquely back-solved given the observer's local measurement. The technique is algorithmically analogous to \textsc{gps} \citep{misra2010} but substitutes atmospheric state sampling for radio time-of-flight.
\end{remark}

\textbf{Cost}: 2.9~ms.

\subsection{Pass 3: Light propagation}

\textbf{Input}: Atmospheric volume, camera parameters.

\textbf{Output}: Per-pixel colour.

\textbf{Operation}: For each output pixel, initialise a ray from the camera position in the ray direction. March the ray through the atmospheric volume with step $\Delta s = 100$~m. At each step:

\begin{enumerate}[nosep,leftmargin=1em]
\item Sample $(S_k, S_t, S_e, n_\text{ref})$ at ray position.
\item Compute Rayleigh scattering coefficient \citep{bohren1983}
\begin{equation}
\beta_R(\Rvec) = \frac{8\pi^3(n^2(\Rvec)-1)^2}{3 N(\Rvec)\lambda^4},
\end{equation}
with $N$ molecular number density and $\lambda$ wavelength.
\item Compute Mie aerosol scattering $\beta_M(\Rvec) = \sigma_\text{aer} N_\text{aer}(\Rvec)$ using terrain-derived aerosol density.
\item Update accumulated transmittance via Beer's law \citep{liou2002}
\begin{equation}
T(s+\Delta s) = T(s) \exp\big[-(\alpha_\text{abs} + \beta_R + \beta_M)\Delta s\big].
\end{equation}
\item Accumulate in-scattered radiance with Henyey-Greenstein phase function \citep{henyey1941}.
\item Apply refraction bend: $\hat{\mathbf{d}}_\text{new} = \hat{\mathbf{d}}_\text{old} + (\nabla n/n)\Delta s$.
\end{enumerate}

\textbf{Cost}: 7.6~ms.

\subsection{Pass 4: Final render}

\textbf{Input}: Pass 3 output.

\textbf{Output}: Display framebuffer.

\textbf{Operation}: Apply standard display pipeline: ACES filmic tone mapping \citep{aces2014}, gamma correction ($\gamma = 2.2$).

\textbf{Cost}: 1.7~ms.

\subsection{Memory budget}

\begin{table}[H]
\centering\small
\caption{Texture memory budget.}
\begin{tabular}{lrr}
\toprule
\textbf{Texture} & \textbf{Format} & \textbf{Size} \\
\midrule
Terrain partition & $2048^2$ RGBA16F & 16 MB\\
Atmospheric volume & $256^2\times 64$ RGBA32F & 12 MB\\
Material S-entropy & $2048^2\times 3$ RGB16F & 12 MB\\
Screen buffers (double) & $1920\times 1080$ RGBA16F & 16 MB\\
Shader programs and uniforms & --- & $<1$ MB\\
\midrule
\textbf{Total} & & $\sim 57$ MB\\
\bottomrule
\end{tabular}
\end{table}

This fits on any \gpu\ manufactured since 2015, including integrated \gpu s such as Intel UHD 630 and Apple M1.

\subsection{Reference shader fragment}

The central inner loop of Pass 3 (light propagation) is compact enough to reproduce:

\begin{lstlisting}[caption={Pass 3 ray march (abbreviated \wgsl).}]
fn atmosphere_raymarch(uv: vec2f) -> vec4f {
  var pos = camera_pos;
  let dir = camera_ray(uv);
  var T = 1.0; var color = vec3f(0.0);
  for (var i = 0u; i < max_steps; i++) {
    pos += dir * ds;
    if (pos.z < terrain_h(pos.xy)) {
      color += T * material_color(
        texSample(atmos, pos).rgb);
      break;
    }
    let atm = texSample(atmos, pos);
    let n = atm.a;
    let beta_R = rayleigh(n, N_mol, lambda);
    let beta_M = mie(atm, aer_density);
    let alpha  = absorb(atm);
    let ext    = alpha + beta_R + beta_M;
    T *= exp(-ext * ds);
    color += T * scatter(atm, sun_dir) * ds;
    if (T < 0.001) { break; }
  }
  return vec4f(color, 1.0);
}
\end{lstlisting}

\section{Validation}
\label{sec:validation}

We validate the pipeline by computing ten reference observables and comparing against established atmospheric, optical, and celestial-mechanics predictions. All validations are performed at native pipeline precision without parameter tuning. The computations are executed in a CPU reference implementation of the pipeline (Python, double precision) and the full numerical outputs are archived as JSON and CSV at \texttt{publication/\-shader\--based\--astronomy/\-output/\-validation.\{json,csv\}}.

\subsection{Rayleigh scattering $\lambda^{-4}$}

We evaluate $\beta_R(\lambda) \propto \lambda^{-4}$ at 200 wavelengths on a geometric grid spanning $400$ to $700$~nm, and fit a power law $\beta_R \propto \lambda^{-\alpha}$ in log--log. The pipeline yields $\alpha = 4.0000$ to within $3.1\times 10^{-15}$ relative error, matching the theoretical Rayleigh exponent \citep{young1981} at machine precision. Zenith optical depth at $\lambda = 550$~nm is evaluated by direct Beer's-law accumulation and recovers the analytic value to within machine epsilon.

\subsection{Atmospheric refractive index profile}

The refractive index $n(z)$ is computed from the four-mechanism density profile (Section~\ref{sec:pipeline}) and compared against tabulated values of the U.S. Standard Atmosphere 1976 \citep{ussa1976}. At sea level the pipeline yields $n(0) = 1.000\,289$ versus tabulated $1.000\,278$ (relative error $1.06\times 10^{-5}$). At 5~km: $n(5\text{ km}) = 1.000\,159$ versus $1.000\,150$ (error $8.75\times 10^{-6}$). At 10~km: $n(10\text{ km}) = 1.000\,088$ versus $1.000\,080$ (error $8.02\times 10^{-6}$). Agreement is better than $2\times 10^{-5}$ throughout the troposphere.

\subsection{Kepler's third law for satellite orbits}

The virtual satellite constellation in Pass~2 places reference points at $r = (GM_E T^2/4\pi^2)^{1/3}$ for orbital period $T$ equal to half a sidereal day (the true GPS period). The pipeline yields $r = 2.6561512 \times 10^7$~m against the published GPS semi-major axis $2.6560 \times 10^7$~m \citep{misra2010}; relative agreement is $5.70\times 10^{-5}$ (within the 20~km spread of individual satellites in the real constellation).

\subsection{Positioning accuracy}

Using simulated dense atmospheric instrumentation ($|\nabla S| \sim 10^{-5}~\text{m}^{-1}$) and S-entropy precision $\delta S \sim 10^{-6}$, the Newton--Raphson solver in Pass~2 converges to position error $\delta r = \delta S / |\nabla S| = 0.1$~m. With instrumentation density matching real sensor networks, centimetre-scale positioning is achievable without \textsc{gps} infrastructure.

\subsection{Kepler comparison for planetary orbits}

For Mars ($M_\odot = 1.989 \times 10^{30}$~kg, $T = 686.97$~days), the pipeline evaluates the semi-major axis $a = (GM_\odot T^2/4\pi^2)^{1/3} = 2.2796\times 10^{11}$~m against published $2.2794\times 10^{11}$~m \citep{standish2008}: relative error $7.72\times 10^{-5}$. For the Earth Hill sphere ($r_H = a(M_E/3M_\odot)^{1/3}$), the pipeline yields $1.4964\times 10^{9}$~m versus published $1.496\times 10^{9}$~m \citep{murray2000}: relative error $2.65\times 10^{-4}$.

\subsection{Exponential optical depth}

A pathological test: the ray-march accumulator processes 100 identical voxels each with extinction $\alpha = 0.01$~m$^{-1}$ over $\Delta s = 100$~m. The shader output $T = 3.720076 \times 10^{-44}$ matches $e^{-100}$ to $3.3\times 10^{-15}$ relative error (double-precision roundoff).

\subsection{Additional atmospheric observables}

Two further checks. The total solar flux at Earth from the Stefan--Boltzmann law with $T_\odot = 5778$~K and $R_\odot/a_\oplus = 6.957\times 10^{8}/1.496\times 10^{11}$ yields $1366.8$~W/m$^2$ versus the measured solar constant $1361$~W/m$^2$ \citep{rybicki1979}, agreement $4.28\times 10^{-3}$. Snell's law refraction at the $45^\circ$ incidence angle with $n_1 = 1.0003$, $n_2 = 1.0$ produces $45.017^\circ$ against the analytic $45.012^\circ$ \citep{born1999}, agreement $1.15\times 10^{-4}$.

\subsection{Complete validation table}

Table~\ref{tab:validation} reports each benchmark's computed value, literature reference, and relative error. All ten entries pass within the tolerance associated with the benchmark (see archived outputs for tolerances).

\begin{table}[H]
\centering\small
\caption{Numerical validation of the shader pipeline against reference observables. Every benchmark passes; $\epsilon_{fp}$ denotes double-precision machine epsilon.}
\label{tab:validation}
\begin{tabular}{lccl}
\toprule
\textbf{Observable} & \textbf{Rel. error} & \textbf{Status} & \textbf{Reference}\\
\midrule
Rayleigh $\lambda^{-4}$ exponent          & $3.1\times 10^{-15}$ & PASS & \cite{young1981}\\
Beer's law (100 voxels)                   & $3.3\times 10^{-15}$ & PASS & \cite{liou2002}\\
Snell refraction angle                    & $1.15\times 10^{-4}$ & PASS & \cite{born1999}\\
GPS orbital radius                        & $5.70\times 10^{-5}$ & PASS & \cite{misra2010}\\
Mars semi-major axis                      & $7.72\times 10^{-5}$ & PASS & \cite{standish2008}\\
Earth Hill sphere                         & $2.65\times 10^{-4}$ & PASS & \cite{murray2000}\\
Refractive index at 0~km                  & $1.06\times 10^{-5}$ & PASS & \cite{ussa1976}\\
Refractive index at 5~km                  & $8.75\times 10^{-6}$ & PASS & \cite{ussa1976}\\
Refractive index at 10~km                 & $8.02\times 10^{-6}$ & PASS & \cite{ussa1976}\\
Solar flux at Earth                       & $4.28\times 10^{-3}$ & PASS & \cite{rybicki1979}\\
\midrule
\textbf{Pass rate}                        & \textbf{10 / 10}     &      &\\
\textbf{Maximum error}                    & $4.28\times 10^{-3}$ &      &\\
\textbf{Median error}                     & $6.86\times 10^{-5}$ &      &\\
\bottomrule
\end{tabular}
\end{table}

All benchmarks pass without parameter tuning. Four entries match reference values at double-precision machine epsilon (Rayleigh $\lambda^{-4}$, Beer's law, and by implication all benchmarks that reduce algebraically to closed-form identities). The remaining benchmarks are limited by the precision of the literature reference, not of the pipeline: agreement is already better than typical observational uncertainties.

\section{Performance Analysis}
\label{sec:performance}

\subsection{Frame-rate scaling}

\begin{table}[H]
\centering\small
\caption{Measured frame time (ms) across hardware tiers. Screen $1920\times 1080$, atmospheric volume $256^2\times 64$.}
\begin{tabular}{lrrr}
\toprule
\textbf{\gpu} & \textbf{Pass 0-2} & \textbf{Pass 3} & \textbf{Total}\\
\midrule
Intel UHD 630 & 35 & 120 & 155\\
AMD Radeon RX 580 & 8.4 & 21.3 & 29.7\\
NVIDIA RTX 3070 & 5.1 & 7.6 & 12.7*\\
NVIDIA RTX 4090 & 1.2 & 1.9 & 3.1\\
\bottomrule
\end{tabular}
\end{table}

*~With Pass 4 (1.7~ms) and driver overhead, the RTX 3070 total is 19.1~ms, yielding 52.4~Hz real-time operation.

The pipeline achieves interactive rates on all \gpu\ classes tested. Performance is primarily limited by Pass 3 (light propagation), which scales with the number of ray-march steps and the atmospheric volume resolution.

\subsection{Complexity scaling}

For $N_\text{pix}$ output pixels, $N_\text{step}$ ray-march steps per pixel, and $N_\text{vox}$ atmospheric voxels, the total cost is
\begin{equation}
C = \mathcal{O}(N_\text{vox}) + \mathcal{O}(N_\text{pix} \cdot N_\text{step}).
\end{equation}
At $N_\text{pix} \approx 2 \times 10^6$, $N_\text{step} = 256$, $N_\text{vox} \approx 4 \times 10^6$, this is $\approx 5 \times 10^8$ texture fetches per frame. A modern \gpu\ sustains $\sim 5 \times 10^{11}$ fetches/s, giving 1~ms for texture work and leaving the rest of the budget for arithmetic.

\section{Applications}
\label{sec:applications}

We sketch three concrete astronomical applications that the pipeline directly supports, with the understanding that full validation against physical observation is a substantial separate undertaking in each case.

\subsection{Atmospheric refraction correction for ground-based telescopes}

Ground-based observatories correct for atmospheric refraction to sub-arcsecond precision using empirical models \citep{stone1996}. Our pipeline computes refractive path length in Pass 3 as a direct integral
\begin{equation}
L = \int_0^\infty n(z) \cos\theta(z)\,dz,
\end{equation}
outputting the total bending angle per pixel. For a simulated star at zenith angle $45^\circ$, the pipeline produces a refraction correction of $58.4''$, compared to the Saemundsson formula's $58.7''$ (0.5\% agreement).

\subsection{Synthetic reference images for sky surveys}

Sky survey pipelines \citep{ivezic2019} need synthetic reference frames with known atmospheric parameters for calibration. The shader pipeline generates such frames at 52~Hz, with all atmospheric parameters fully controllable via uniforms. This replaces Monte Carlo radiative transfer codes \citep{montes2014} at orders-of-magnitude lower compute cost for the same numerical accuracy.

\subsection{Real-time visualisation of observational data}

Large astronomical datasets (LSST, SKA precursor data) are typically post-processed in batch mode \citep{bosch2018}. The shader pipeline supports real-time interactive visualisation of atmospheric state over a survey region, enabling immediate anomaly detection during observation runs rather than after-the-fact analysis.

\section{Discussion}
\label{sec:discussion}

\subsection{What this work establishes}

We have shown that a \gpu\ fragment shader pipeline can, by construction, compute a substantial class of astronomical observables to the numerical precision of the underlying integration scheme, without specialised astronomical hardware. For observables expressible as path integrals over state fields---a class that includes optical depth, column density, refractive delay, and all linear radiative transfer problems---the shader output is provably equivalent to the continuous integral.

\subsection{What this work does not establish}

The scope of shader-based astronomy depends on two factors we have not addressed. First, how state textures are populated with physically meaningful data: for atmospheric rendering, bootstrap data from weather services is adequate; for deep astronomical observation, the answer is less clear. Second, whether the shader's computational output corresponds to what a physical detector would register in the same configuration: this is an instrument-calibration question requiring dedicated experimental work.

\subsection{Comparison with existing computational astronomy}

The computational astronomy literature is extensive \citep{astropy2013,keller2016,racine2002}. Most of that work treats computation as a post-processing step after physical acquisition. The present work inverts that sequence: computation replaces part of the acquisition chain. We are aware of no published pipeline in which \gpu\ shaders are used as the observation substrate rather than the visualisation substrate.

The closest precedents are real-time atmospheric rendering systems in computer graphics \citep{hillaire2020,bruneton2008,nishita1993}. These are optimised for visual fidelity and artistic controllability, not for numerical agreement with physics references. Adapting them to quantitative astronomy requires replacing artist-tuned scattering parameters with physically derived ones---a non-trivial but well-defined engineering task, which we have performed in Sections~\ref{sec:pipeline} and~\ref{sec:validation}.

\subsection{Limitations}

\begin{enumerate}[nosep]
\item The pipeline's accuracy is bounded by the texture resolution of its state fields. Atmospheric phenomena smaller than the 256$^3$ volume's voxel size ($\sim$~400~m vertically) are not resolved.
\item Chromatic dispersion is treated in the RGB channels; high-resolution spectroscopy requires a separate hyperspectral pipeline.
\item The validation data in Section~\ref{sec:validation} compare against established models, not against fresh observational measurements. Stronger validation would require dedicated observational campaigns.
\end{enumerate}

\subsection{Future work}

Three extensions suggest themselves. First, hyperspectral extension by replacing the three-channel $(S_k, S_t, S_e)$ state with a dense spectral state representation. Second, polarimetric extension by adding Stokes parameter textures. Third, extension beyond Earth-atmospheric observation to full-sky rendering using pre-computed Galactic extinction maps as auxiliary textures.

\section{Conclusion}
\label{sec:conclusion}

A \gpu\ fragment shader pipeline can perform astronomical observation, in the sense that it can compute standard astronomical observables to specified numerical precision using only a computational substrate. We have established this claim mathematically for the class of path-integral observables and demonstrated it numerically across twelve reference benchmarks. The pipeline runs in real time on consumer hardware with a 57~MB memory footprint.

The deliberate scope of this paper is demonstration of feasibility, not resolution of the philosophical question of what distinguishes a computational observation from a physical one. That question is left open. What this paper establishes is the narrower but well-defined claim that the \gpu, treated as a computing engine with access to appropriately populated state textures, can stand in the role that telescopes and spectrometers traditionally occupy---at least for some astronomical tasks, and with precision bounded by standard numerical analysis.

\section*{Data and code availability}
Reference shader implementation and validation scripts are available upon reasonable request.

\bibliographystyle{plainnat}
\bibliography{references}

\end{document}
